% ============================================================
%  Basado en la plantilla de la Practica 5 (mismos colores/estilos)
% ============================================================
\documentclass[11pt,a4paper]{article}

% ---------- Idioma y codificación ----------
\usepackage[utf8]{inputenc}
\usepackage[T1]{fontenc}
\usepackage[spanish,es-noshorthands]{babel}

% ---------- Tipografía y layout ----------
\usepackage{lmodern}
\usepackage[sfdefault,light]{roboto}
\usepackage{geometry}
\geometry{a4paper, top=3.1cm, bottom=2.8cm, left=2.4cm, right=2.4cm, headheight=1.6cm}
\usepackage{ragged2e}
\usepackage{enumitem}
\usepackage{graphicx}
\usepackage{float}
\usepackage{booktabs}
\usepackage{array}
\usepackage[most]{tcolorbox}
\usepackage{tikz}
\usetikzlibrary{shadows,calc}
\usepackage{titlesec}
\usepackage{fancyhdr}
\usepackage{fontawesome5}
\usepackage[colorlinks=true,breaklinks=true]{hyperref}
\usepackage{caption}
\usepackage{microtype}
\usepackage{xurl}
\usepackage{listings}

% ---------- Paleta de colores (tema "coolnight" de tu terminal Kitty) ----------
\definecolor{cNavy}{HTML}{011423}
\definecolor{cTealDeep}{HTML}{064B5A}
\definecolor{cTeal}{HTML}{097289}
\definecolor{cBlue}{HTML}{0FC5ED}
\definecolor{cCyan}{HTML}{0FC5ED}
\definecolor{cMint}{HTML}{44FFB1}
\definecolor{cPurple}{HTML}{A277FF}
\definecolor{cRed}{HTML}{E52E2E}
\definecolor{cBg}{HTML}{EAF6F8}
\definecolor{cText}{HTML}{01111E}
\definecolor{cOutBg}{HTML}{F4F4F4}

\hypersetup{
  linkcolor=cTealDeep,
  citecolor=cTealDeep,
  urlcolor=cBlue,
  pdftitle={Práctica 6 -- Reglas de Asociación en R},
  pdfauthor={Miguel Deza}
}

\color{cText}
\emergencystretch=2.5em

% ---------- Encabezado / pie de página ----------
\pagestyle{fancy}
\fancyhf{}
\renewcommand{\headrulewidth}{0pt}
\renewcommand{\footrulewidth}{0pt}
\fancyhead[L]{\small\color{cTealDeep}\textbf{\faGraduationCap\ Maestría en Ciencia de Datos}}
\fancyhead[R]{\small\color{cTeal}Práctica 6 -- Reglas en R}
\fancyfoot[C]{%
  \begin{tikzpicture}[remember picture]
    \node[circle, fill=cTealDeep, minimum size=1.7em, inner sep=0pt] (pg) {\color{white}\small\thepage};
  \end{tikzpicture}%
}
\renewcommand{\headrule}{{\color{cCyan}\hrule height 1.4pt}\vspace{1pt}{\color{cBlue}\hrule height 0.5pt}}

% ---------- Estilo de secciones ----------
\titleformat{\section}
  {\normalfont\Large\bfseries\color{cTealDeep}}
  {\tikz[baseline=-0.6ex]{\node[fill=cCyan, text=white, rounded corners=1.5pt, inner sep=3pt, font=\bfseries]{\thesection};}}
  {0.6em}{}
  [{\vspace{-0.3em}\color{cCyan}\titlerule[1.2pt]\vspace{0.3em}}]

\titleformat{\subsection}
  {\normalfont\large\bfseries\color{cTeal}}
  {\color{cTeal}\thesubsection}{0.6em}{}

\titlespacing*{\section}{0pt}{1.4em}{0.8em}
\titlespacing*{\subsection}{0pt}{1.1em}{0.5em}

% ---------- Figura sin marco ----------
\newcommand{\figuraMarcada}[3]{%
  \begin{figure}[H]
    \centering
    \includegraphics[width=#2\linewidth]{#1}
    \caption{#3}
  \end{figure}
}

\captionsetup{labelfont={bf,color=cTealDeep}, font=small, textfont=it}

% ---------- Bloques de codigo R y salida de consola ----------
\lstdefinestyle{styleR}{
  language=R,
  backgroundcolor=\color{cBg},
  basicstyle=\ttfamily\small\color{cText},
  keywordstyle=\color{cTealDeep}\bfseries,
  commentstyle=\color{gray}\itshape,
  stringstyle=\color{cTeal},
  numbers=none,
  breaklines=true,
  frame=l,
  framerule=2.2pt,
  rulecolor=\color{cCyan},
  xleftmargin=8pt,
  framexleftmargin=6pt,
  aboveskip=6pt, belowskip=6pt,
  showstringspaces=false
}

\lstdefinestyle{styleOut}{
  backgroundcolor=\color{cOutBg},
  basicstyle=\ttfamily\scriptsize\color{cText},
  breaklines=true,
  frame=l,
  framerule=2.2pt,
  rulecolor=\color{cTealDeep},
  xleftmargin=8pt,
  framexleftmargin=6pt,
  aboveskip=2pt, belowskip=6pt,
  showstringspaces=false
}

\newcommand{\codigoR}[1]{\lstinputlisting[style=styleR]{#1}}

% ============================================================
\begin{document}

% -------------------- PORTADA --------------------
\begin{titlepage}
\newgeometry{margin=0cm}
\begin{tikzpicture}[remember picture, overlay]
  \fill[cNavy] (current page.south west) rectangle (current page.north east);

  \fill[cPurple, opacity=0.28] (current page.north east) ++(-1,1) circle (3.2);
  \fill[cCyan, opacity=0.22] (current page.north east) ++(-2.6,0.4) circle (2.1);
  \fill[cTealDeep, opacity=0.55] (current page.south west) ++(1,-1) circle (4.0);
  \fill[cMint, opacity=0.16] (current page.south west) ++(2.4,-0.2) circle (2.3);

  % ---- Icono insignia ----
  \node[circle, fill=cTealDeep, draw=cCyan, line width=1.3pt, minimum size=2.5cm] (logo) at ($(current page.north)+(0,-2.9)$) {};
  \node at (logo) {\color{cMint}\fontsize{28}{28}\selectfont\faIcon{r-project}};

  \node[anchor=north, text width=14cm, align=center] at ($(current page.north)+(0,-4.7)$) {
    {\color{cCyan}\bfseries\fontsize{12}{14}\selectfont \faDatabase\ \ MAESTRÍA EN CIENCIA DE DATOS}
  };

  \node[anchor=center, text width=15cm, align=center] at ($(current page.center)+(0,1.2)$) {
    {\color{white}\bfseries\fontsize{40}{44}\selectfont Práctica 6}\\[0.35em]
    {\color{cMint}\bfseries\fontsize{20}{24}\selectfont \faIcon{r-project}\ Reglas de Asociación en R}
  };

  \draw[cCyan, line width=1.4pt] ($(current page.center)+(-3.6,-0.55)$) -- ++(7.2,0);

  \node[anchor=center, text width=13cm, align=center] at ($(current page.center)+(0,-1.3)$) {
    {\color{white!85!black}\fontsize{11}{14}\selectfont
     Minería de reglas de asociación con el paquete \textbf{arules} en R:\\
     carga, EDA, Apriori y visualización (minsup $\geq$ 0.5, minconf $\geq$ 0.6)}
  };

  \node[anchor=south, text width=14cm, align=center] at ($(current page.south)+(0,2.6)$) {
    \begin{tabular}{r l}
      \color{cCyan}\bfseries \faUser\ Elaborado por: & \color{white} Miguel Deza \\[3pt]
      \color{cCyan}\bfseries \faGraduationCap\ Programa: & \color{white} Maestría en Ciencia de Datos \\[3pt]
      \color{cCyan}\bfseries \faBook\ Curso: & \color{white} Introducción a Ciencia de Datos \\[3pt]
      \color{cCyan}\bfseries \faChalkboardTeacher\ Docente: & \color{white} Yessenia Deysi Yari Ramos \\[3pt]
      \color{cCyan}\bfseries \faIcon{calendar-alt}\ Fecha: & \color{white} 19 de septiembre de 2026 \\
    \end{tabular}
  };
\end{tikzpicture}
\end{titlepage}
\restoregeometry

% -------------------- 1. DESARROLLO --------------------
\section{\faCogs\ \ Desarrollo}

\subsection{\faBoxOpen\ \ 1. Cargar las bibliotecas necesarias}
\begin{lstlisting}[style=styleR]
library(arules)
library(arulesViz)
\end{lstlisting}
Ambos paquetes cargaron sin errores (\texttt{arules} para minar itemsets/reglas, \texttt{arulesViz} para graficarlas).

\subsection{\faFileImport\ \ 2. Cargar la data (CSV)}
\begin{lstlisting}[style=styleR]
datos <- read.csv("dataset_reglas_asociacion.csv", stringsAsFactors = TRUE)
print(datos)
\end{lstlisting}
\begin{lstlisting}[style=styleOut]
  Transaccion  A  B  C  E
1           1 SI SI SI NO
2           2 NO SI SI SI
3           3 SI SI SI SI
4           4 NO SI NO SI
\end{lstlisting}

\subsection{\faSearch\ \ 3. Análisis exploratorio de datos}
\begin{lstlisting}[style=styleR]
str(datos)
summary(datos)
sapply(datos[, c("A", "B", "C", "E")], table)
\end{lstlisting}
\begin{lstlisting}[style=styleOut]
'data.frame': 4 obs. of  5 variables:
 $ Transaccion: int  1 2 3 4
 $ A: Factor w/ 2 levels "NO","SI": 2 1 2 1
 $ B: Factor w/ 1 level "SI": 1 1 1 1
 $ C: Factor w/ 2 levels "NO","SI": 2 2 2 1
 $ E: Factor w/ 2 levels "NO","SI": 1 2 2 2

$A
NO SI
 2  2
$B
SI
 4
$C
NO SI
 1  3
$E
NO SI
 1  3
\end{lstlisting}
El atributo \textbf{B} tiene un solo nivel (siempre ``SI''): no varía entre transacciones. \textbf{A} está balanceado 50/50, mientras que \textbf{C} y \textbf{E} son ``SI'' en el 75\% de los casos.

\subsection{\faTools\ \ 4. Preprocesar los datos}
Se descarta el ID de transacción (no es un ítem) y se convierte SI/NO a lógico -- \texttt{TRUE} indica que el ítem está presente -- para poder construir el objeto \texttt{transactions} que usa \texttt{arules}.
\begin{lstlisting}[style=styleR]
datos_log <- datos[, c("A", "B", "C", "E")]
datos_log <- as.data.frame(lapply(datos_log, function(x) x == "SI"))
transacciones <- as(datos_log, "transactions")
summary(transacciones)
inspect(transacciones)
\end{lstlisting}
\begin{lstlisting}[style=styleOut]
transactions as itemMatrix in sparse format with
 4 rows (elements/itemsets/transactions) and
 4 columns (items) and a density of 0.75

most frequent items:
  B   C   E   A (Other)
  4   3   3   2       0

    items        transactionID
[1] {A, B, C}    1
[2] {B, C, E}    2
[3] {A, B, C, E} 3
[4] {B, E}       4
\end{lstlisting}
\figuraMarcada{img/frecuencia-items.png}{0.75}{Frecuencia relativa (soporte individual) de cada ítem.}

\subsection{\faMagic\ \ 5. Generar las reglas de asociación}
\begin{lstlisting}[style=styleR]
reglas <- apriori(transacciones,
                   parameter = list(supp = 0.5, conf = 0.6, target = "rules"))
inspect(sort(reglas, by = "confidence"))

itemsets <- apriori(transacciones,
                     parameter = list(supp = 0.5, target = "frequent itemsets"))
inspect(sort(itemsets, by = "support"))
\end{lstlisting}
\begin{lstlisting}[style=styleOut]
     lhs       rhs support confidence coverage lift      count
[1]  {}     => {B} 1.00    1.0000000  1.00     1.0000000 4
[2]  {A}    => {C} 0.50    1.0000000  0.50     1.3333333 2
[3]  {A}    => {B} 0.50    1.0000000  0.50     1.0000000 2
[4]  {E}    => {B} 0.75    1.0000000  0.75     1.0000000 3
[5]  {C}    => {B} 0.75    1.0000000  0.75     1.0000000 3
[6]  {A, C} => {B} 0.50    1.0000000  0.50     1.0000000 2
[7]  {A, B} => {C} 0.50    1.0000000  0.50     1.3333333 2
[8]  {C, E} => {B} 0.50    1.0000000  0.50     1.0000000 2
[9]  {}     => {E} 0.75    0.7500000  1.00     1.0000000 3
[10] {}     => {C} 0.75    0.7500000  1.00     1.0000000 3
[11] {B}    => {E} 0.75    0.7500000  1.00     1.0000000 3
[12] {B}    => {C} 0.75    0.7500000  1.00     1.0000000 3
[13] {C}    => {A} 0.50    0.6666667  0.75     1.3333333 2
[14] {E}    => {C} 0.50    0.6666667  0.75     0.8888889 2
[15] {C}    => {E} 0.50    0.6666667  0.75     0.8888889 2
[16] {B, C} => {A} 0.50    0.6666667  0.75     1.3333333 2
[17] {B, E} => {C} 0.50    0.6666667  0.75     0.8888889 2
[18] {B, C} => {E} 0.50    0.6666667  0.75     0.8888889 2

     items     support count
[1]  {B}       1.00    4
[2]  {E}       0.75    3
[3]  {C}       0.75    3
[4]  {B, E}    0.75    3
[5]  {B, C}    0.75    3
[6]  {A}       0.50    2
[7]  {A, C}    0.50    2
[8]  {A, B}    0.50    2
[9]  {C, E}    0.50    2
[10] {A, B, C} 0.50    2
[11] {B, C, E} 0.50    2
\end{lstlisting}
Con \texttt{supp $\geq$ 0.5} y \texttt{conf $\geq$ 0.6}, \texttt{arules} encontró \textbf{11 itemsets frecuentes} y \textbf{18 reglas}. A diferencia del conteo manual, el algoritmo de \texttt{arules} genera por defecto solo reglas con \textbf{un ítem en el consecuente} (por eso no aparecen reglas del tipo A $\to$ \{B,C\}) y sí incluye 3 reglas con antecedente vacío (\texttt{\{\} => \{B\}}, etc.), que representan simplemente el soporte individual de cada ítem frecuente.

\subsection{\faChartBar\ \ 6. Visualizar las reglas generadas}
\begin{lstlisting}[style=styleR]
plot(reglas, measure = c("support", "confidence"), shading = "lift")
plot(reglas, method = "graph", engine = "igraph")
plot(reglas, method = "grouped")
\end{lstlisting}

\figuraMarcada{img/reglas-scatter.png}{0.75}{Soporte vs. confianza de las 18 reglas, coloreadas por lift.}

\figuraMarcada{img/reglas-grafo.png}{0.85}{Grafo de reglas: el tamaño del nodo es el soporte, el color el lift.}

\figuraMarcada{img/reglas-agrupadas.png}{0.85}{Matriz agrupada: reglas agrupadas por antecedente (columnas) y consecuente (filas).}

\end{document}
